\section{Test-Cost-Sensitive Attribute Reduction}
\label{sec:formal-method}
This section formulates the test-cost-sensitive attribute reduction problem for label distribution learning (LDL) as a constrained combinatorial optimization task. In contrast to classical rough-set frameworks that enforce dependency preservation as a rigid feasibility constraint, our formulation employs an inner validation loss threshold to govern predictive feasibility, while leveraging the proposed granular dependency $\Gamma(P)$ to guide the search trajectory toward cost-effective subsets.

\subsection{Prediction-Feasible Minimum-Cost Formulation}
Let $c_a > 0$ denote the acquisition or measurement cost incurred when querying attribute $a \in A$ at inference time, and let $C(P) = \sum_{a \in P} c_a$ represent the total acquisition cost of subset $P \subseteq A$, with $C(\varnothing) = 0$.

To ensure that attribute reduction does not compromise predictive fidelity, feasibility is defined over an inner cross-validation protocol on the training set $O$. Let $O$ be partitioned into $K$ disjoint validation folds $I_1, \dots, I_K$. For an evaluated subset $P \subseteq A$, let $\widehat{\boldsymbol{\ell}}^{(-k)}_P(\cdot)$ denote the LDL model trained on $O \setminus I_k$ using features restricted to $P$. The inner validation loss is given by
\begin{equation}
	R_{\mathrm{inner}}(P) = \frac{1}{n} \sum_{k=1}^K \sum_{i \in I_k} \|\boldsymbol{\ell}_i - \widehat{\boldsymbol{\ell}}^{(-k)}_P(x_i)\|_2^2.
	\label{eq:inner-loss}
\end{equation}
We define the feasibility baseline as the performance attained by the complete attribute set, $L_{\max} = R_{\mathrm{inner}}(A)$. A candidate subset $P$ is deemed \emph{prediction-feasible} if its predictive loss does not exceed this baseline:
\begin{equation}
	\mathcal{F} = \left\{ P \subseteq A : R_{\mathrm{inner}}(P) \le L_{\max} \right\}.
	\label{eq:feasible-set}
\end{equation}
The optimal test-cost-sensitive reduction problem is then formalized as:
\begin{equation}
	P^* \in \underset{P \in \mathcal{F}}{\operatorname{argmin}} \; C(P).
	\label{eq:formal-cost-objective}
\end{equation}

By construction, the full set satisfies feasibility ($A \in \mathcal{F}$). If the baseline prediction of the empty set satisfies $R_{\mathrm{inner}}(\varnothing) \le L_{\max}$, then $P^* = \varnothing$ is the unique zero-cost trivial solution. In our reference implementation, $\widehat{\boldsymbol{\ell}}^{(-k)}_P$ is instantiated as a standard 5-nearest-neighbor regressor operating on min--max scaled attributes with Euclidean distance, while the empty subset predicts the fold-wise mean distribution $\bar{\boldsymbol{\ell}}^{(-k)}$.

\begin{remark}[Decoupling Feasibility from Dependency]
	Classical rough-set reducts typically mandate $\Gamma(P) \ge \Gamma(A)$. However, because empirical label dispersion and downstream non-parametric regression losses do not induce strictly isomorphic partial orders over $2^A$, imposing $\Gamma(P) \ge \Gamma(A)$ can unnecessarily discard subsets with superior test-cost trade-offs. Eq.~\eqref{eq:formal-cost-objective} instead designates actual predictive fidelity ($R_{\mathrm{inner}}$) as the feasibility boundary and reserves $\Gamma(P)$ as an efficient heuristic surrogate during search.
\end{remark}

\subsection{Dependency-Guided Constructive Search}
\label{sub-sec:dependency}
Finding $P^*$ in Eq.~\eqref{eq:formal-cost-objective} is NP-hard. We design a multi-tier forward-selection and backward-elimination heuristic that prioritizes predictive gain per unit cost, gated by granular dependency.

For any subset $P$ and an unselected candidate $a \in A \setminus P$, define the marginal predictive gain $u_a(P)$ and marginal granular dependency gain $v_a(P)$ as:
\begin{equation}
	u_a(P) = R_{\mathrm{inner}}(P) - R_{\mathrm{inner}}(P \cup \{a\}), \qquad
	v_a(P) = \Gamma(P \cup \{a\}) - \Gamma(P).
	\label{eq:marginal-gains}
\end{equation}

\subsubsection{Forward Selection Stage}
Starting from $P = \varnothing$, attributes are greedily accumulated until $P \in \mathcal{F}$. At each iteration, candidate attributes in $A \setminus P$ are partitioned hierarchically:
\begin{enumerate}[label=(\roman*)]
	\item \textbf{Loss-improving pool:} $J(P) = \{a \in A \setminus P : u_a(P) > 0\}$.
	\item \textbf{Dual-improving pool:} $H(P) = \{a \in J(P) : v_a(P) > 0\}$.
\end{enumerate}
The next attribute is selected according to the following priority hierarchy:
\begin{equation}
	a^* = 
	\begin{cases}
		\underset{a \in H(P)}{\operatorname{argmax}} \; \dfrac{v_a(P)}{c_a}, & \text{if } H(P) \ne \varnothing, \\[8pt]
		\underset{a \in J(P)}{\operatorname{argmax}} \; \dfrac{u_a(P)}{c_a}, & \text{if } H(P) = \varnothing \text{ and } J(P) \ne \varnothing, \\[8pt]
		\underset{a \in A \setminus P}{\operatorname{argmax}} \; u_a(P), & \text{otherwise (plateau-bridging step)},
	\end{cases}
	\label{eq:selection-rule}
\end{equation}
with ties resolved by favoring smaller acquisition costs $c_a$, followed by fixed attribute indices. The plateau-bridging branch allows the search to navigate non-convex valleys where no single attribute produces an immediate positive marginal reduction in loss.

\subsubsection{Restarting Backward Elimination}
Upon attaining feasibility ($P \in \mathcal{F}$), redundant attributes are pruned. Attributes in $P$ are sorted in descending order of cost $c_a$ (breaking ties by ascending index). For each attribute $a \in P$:
\begin{itemize}
	\item If $R_{\mathrm{inner}}(P \setminus \{a\}) \le L_{\max}$, the removal is committed ($P \gets P \setminus \{a\}$), and the pruning scan restarts from the highest-cost attribute of the updated set.
	\item The elimination terminates when a full scan over all attributes in $P$ yields no feasible deletions.
\end{itemize}
The complete procedure is formalized in Algorithm~\ref{alg2:guided}.

\begin{algorithm}[t]
	\caption{Dependency-Guided Cost-Sensitive Attribute Reduction}
	\label{alg2:guided}
	\begin{algorithmic}[1]
		\Require Attribute set $A$, acquisition costs $\{c_a\}_{a \in A}$, baseline loss $L_{\max} = R_{\mathrm{inner}}(A)$
		\Ensure Feasible, 1-deletion irreducible attribute subset $P^*$
		\State Initialize $P \gets \varnothing$
		\If{$R_{\mathrm{inner}}(\varnothing) \le L_{\max}$}
		\State \Return $\varnothing$
		\EndIf
		\While{$R_{\mathrm{inner}}(P) > L_{\max}$ \textbf{and} $P \ne A$} \Comment{Forward Selection}
		\State Compute $u_a(P)$ and $v_a(P)$ for all $a \in A \setminus P$ via Eq.~\eqref{eq:marginal-gains}
		\State Form $J(P) = \{a : u_a(P) > 0\}$ and $H(P) = \{a \in J(P) : v_a(P) > 0\}$
		\State Select $a^*$ according to the tiered priority in Eq.~\eqref{eq:selection-rule}
		\State $P \gets P \cup \{a^*\}$
		\EndWhile
		\State \textbf{repeat} \Comment{Restarting Backward Pruning}
		\State Sort attributes in $P$ such that $c_{a_1} \ge c_{a_2} \ge \dots \ge c_{a_{|P|}}$
		\State $\mathrm{pruned} \gets \textbf{false}$
		\For{$j = 1$ \textbf{to} $|P|$}
		\If{$R_{\mathrm{inner}}(P \setminus \{a_j\}) \le L_{\max}$}
		\State $P \gets P \setminus \{a_j\}$; $\mathrm{pruned} \gets \textbf{true}$; \textbf{break}
		\EndIf
		\EndFor
		\State \textbf{until} $\mathrm{pruned} = \textbf{false}$
		\State \Return $P^* \gets P$
	\end{algorithmic}
\end{algorithm}



\begin{proposition}[Termination and Irreducibility]
	\label{prop:heuristic-guarantee}
	Under deterministic subset evaluations, Algorithm~\ref{alg2:guided} terminates in at most $m$ additions and $m$ accepted deletions, where $m = |A|$. The resulting subset $P^*$ is guaranteed to be feasible ($P^* \in \mathcal{F}$) and single-deletion irreducible:
	\begin{equation}
		\forall a \in P^*, \quad R_{\mathrm{inner}}(P^* \setminus \{a\}) > L_{\max}.
	\end{equation}
\end{proposition}
\begin{proof}
	Each forward step adds an unused attribute from $A \setminus P$. In the worst-case scenario where intermediate subsets remain infeasible, the forward phase expands to $P = A$, which satisfies $R_{\mathrm{inner}}(A) \le L_{\max}$ by definition. Thus, the forward phase terminates in at most $m$ iterations. 
	
	In the backward phase, each accepted deletion strictly decreases $|P|$ by 1 while maintaining $R_{\mathrm{inner}}(P) \le L_{\max}$. Hence, at most $m$ deletions can occur. The phase halts only after a complete scan over all members in $P$ confirms that no single element can be excised without exceeding $L_{\max}$, establishing single-deletion irreducibility.
\end{proof}

\subsection{Cost-Reducing Exchange Refinement}
\label{sub-sec:optional}
To escape sub-optimal local minima reached by forward-backward passes, we incorporate an optional 1-exchange neighborhood search over the feasible space.

Given a feasible subset $P \in \mathcal{F}$, define the set of cost-improving replacement candidates:
\begin{equation}
	\mathcal{N}(P) = \left\{ Q = (P \setminus \{a\}) \cup \{b\} : a \in P, \, b \in A \setminus P, \, c_b < c_a \right\}.
\end{equation}
Candidates in $\mathcal{N}(P)$ are prioritized by decreasing cost saving $\Delta C = c_a - c_b$. Within a predefined computational evaluation budget $B_{\mathrm{eval}}$ (e.g., $B_{\mathrm{eval}} = 200$ subset queries):
\begin{enumerate}
	\item Candidates are sequentially evaluated. Any candidate satisfying $R_{\mathrm{inner}}(Q) \le L_{\max}$ is flagged as admissible.
	\item If admissible exchanges exist, the move yielding the minimal total cost $C(Q)$ (breaking ties by lower inner loss $R_{\mathrm{inner}}(Q)$) is committed: $P \gets Q$.
	\item After committing an exchange, backward pruning is re-invoked to eliminate any secondary redundancies exposed by the new feature combination.
\end{enumerate}
Because each accepted move strictly lowers the acquisition cost ($C(Q) < C(P)$) while preserving feasibility, the sequence of solutions is strictly cost-decreasing, guaranteeing convergence without cycling.

\subsection{Exact Benchmark Formulation and Complexity Characterization}
\label{subsec:exact-search}
To establish a rigorous ground truth for algorithmic benchmarking on moderate-dimensional datasets, Eq.~\eqref{eq:formal-cost-objective} can be solved globally via a cost-bounded Branch-and-Bound (B\&B) search on the subset lattice $2^A$:
\begin{itemize}
	\item \textbf{Incumbent initialization:} The incumbent solution is initialized to the full set, $P_{\mathrm{best}} \gets A$, with incumbent cost $C^* = C(A)$.
	\item \textbf{Branching and Pruning:} The search space is explored via a binary inclusion/exclusion tree. For any partial assignment with committed attributes $P_{\mathrm{curr}}$, the node is immediately pruned if its lower-bound cost satisfies:
	\begin{equation}
		C(P_{\mathrm{curr}}) \ge C^*.
		\label{eq:cost-prune}
	\end{equation}
	Since attribute costs are strictly positive ($c_a > 0$), all completions $Q \supseteq P_{\mathrm{curr}}$ satisfy $C(Q) > C(P_{\mathrm{curr}}) \ge C^*$, rendering further branching redundant.
	\item \textbf{Feasibility Update:} When a candidate subset $P$ with $C(P) < C^*$ satisfies $R_{\mathrm{inner}}(P) \le L_{\max}$, the incumbent is updated: $P_{\mathrm{best}} \gets P$, $C^* \gets C(P)$.
\end{itemize}

\subsubsection{Computational Complexity}
The computational characteristics of the reduction pipeline are summarized as follows:
\begin{itemize}
	\item \textbf{Heuristic Search:} The forward phase requests at most $\sum_{k=1}^m (m - k + 1) = O(m^2)$ subset evaluations. The restarting backward elimination performs at most $O(m^2)$ loss evaluations in the worst case. Thus, the total number of subset loss requests is bounded by $O(m^2)$.
	\item \textbf{Exchange Refinement:} Hard-bounded by $\min\{|\mathcal{N}(P)|, B_{\mathrm{eval}}\}$, ensuring controlled wall-clock execution.
	\item \textbf{Exact Search:} Requires $O(2^m)$ evaluations in the theoretical worst case, though cost-based pruning significantly contracts the practical search space when low-cost feasible subsets exist early in the traversal.
\end{itemize}

\begin{remark}[Scope and Generalization Integrity]
	All metric projections, feature scalings, and neighborhood models are strictly fit on the inner training folds $O \setminus I_k$, preventing information leakage. While membership in $\mathcal{F}$ certifies internal feasibility against the full-feature baseline on the development set, downstream performance preservation is evaluated externally on completely isolated hold-out test folds, as detailed in Section~\ref{sec:paper-experiments}.
\end{remark}

\begin{algorithm}[htbp]
	\caption{Dependency-Guided Prediction-Feasible Reduction}
	\label{alg3:guided}
	\begin{algorithmic}[1]
		\Require Attribute set $A=\{0,\ldots,m-1\}$, attribute costs $\{c_a > 0\}_{a \in A}$, loss evaluation oracle $R_{\mathrm{inner}}(\cdot)$, granular dependency oracle $\Gamma(\cdot)$.
		\Ensure Prediction-feasible subset $P^*$ and its acquisition cost $C(P^*)$.
		\State Initialize $L_{\max} \gets R_{\mathrm{inner}}(A)$, $P \gets \varnothing$
		\If{$R_{\mathrm{inner}}(\varnothing) \le L_{\max}$}
		\State \Return $\varnothing, 0$ \Comment{Empty set is trivially feasible and optimal}
		\EndIf
		\While{$R_{\mathrm{inner}}(P) > L_{\max}$} \Comment{Forward selection phase}
		\State $U \gets A \setminus P$, \quad $r \gets R_{\mathrm{inner}}(P)$
		\State Compute $u_a \gets r - R_{\mathrm{inner}}(P \cup \{a\})$ for all $a \in U$
		\State $J \gets \{a \in U : u_a > 0\}$
		\If{$J \ne \varnothing$}
		\State Compute $v_a \gets \Gamma(P \cup \{a\}) - \Gamma(P)$ for all $a \in J$
		\State $H \gets \{a \in J : v_a > 0\}$
		\Else
		\State $H \gets \varnothing$
		\EndIf
		\If{$H \ne \varnothing$}
		\State $a^* \gets \operatorname{argmax}_{a \in H} (v_a / c_a)$ \Comment{Dual-gain guided step}
		\ElsIf{$J \ne \varnothing$}
		\State $a^* \gets \operatorname{argmax}_{a \in J} (u_a / c_a)$ \Comment{Prediction-gain step}
		\Else
		\State $a^* \gets \operatorname{argmax}_{a \in U} u_a$ \Comment{Plateau-bridging step}
		\EndIf
		\State \textbf{break ties} by smaller $c_a$, then smaller index $a$
		\State $P \gets P \cup \{a^*\}$
		\EndWhile
		\Repeat \Comment{Restarting backward pruning phase}
		\State Sort $P$ in descending order of cost $c_a$ (ties broken by ascending index)
		\State $\mathrm{pruned} \gets \textbf{false}$
		\For{$a \in P$}
		\If{$R_{\mathrm{inner}}(P \setminus \{a\}) \le L_{\max}$}
		\State $P \gets P \setminus \{a\}$, \quad $\mathrm{pruned} \gets \textbf{true}$
		\State \textbf{break} \Comment{Restart backward scan immediately}
		\EndIf
		\EndFor
		\Until{$\mathrm{pruned} = \textbf{false}$}
		\State \Return $P^* \gets P, \; C(P^*) \gets \sum_{a \in P^*} c_a$
	\end{algorithmic}
\end{algorithm}
\FloatBarrier
\subsection{Computational Complexity and Scalability Analysis}
\label{subsec:complexity}
Let $n$ denote the number of training instances, $m = |A|$ the total attribute cardinality, $q$ the label dimension, $K$ the number of inner validation folds, and $k$ the neighbor cardinality in the base predictor. We analyze the asymptotic computational complexity and memory consumption of the proposed framework, explicitly distinguishing offline preprocessing from online combinatorial search.

\subsubsection{Preprocessing and Single-Evaluation Oracles}
\begin{enumerate}[label=(\roman*)]
	\item \textbf{Geometric Statistics and Cache Preprocessing ($T_{\mathrm{prep}}$):}
	Precomputing dense pairwise affinity matrices and kernel bandwidths for each attribute across $n$ instances requires $O(m n^2 \log n)$ operations using comparison-based sorting (or $O(m n^2)$ via linear-time median selection). Multiplying affinity matrices by the $n \times q$ label distribution matrix to derive local label profiles costs $O(m n^2 q)$. Storing the precomputed per-attribute distance matrices, instance labels, and feature caches requires $M_{\mathrm{prep}} = O(m n^2 + n q + n m)$ spatial memory.
	
	\item \textbf{Inner Prediction-Loss Evaluation ($T_R(s)$):}
	Evaluating a candidate subset $P$ of size $s = |P|$ requires aggregating the precomputed pairwise distance components across the $K$ validation folds. Partitioning the training instances into validation and fitting sets across $K$ folds yields an aggregation cost of $O(n^2 s)$. Identifying the $k$-nearest neighbors via partial sorting across the folds and aggregating neighbor distributions incurs $O(n^2 \log n + n k q)$ operations. Hence, a single uncached loss evaluation satisfies:
	\begin{equation}
		T_R(s) = O\left(n^2(s + \log n) + n k q\right).
		\label{eq:eval-loss-complexity}
	\end{equation}
	The baseline empty set $P = \varnothing$ bypasses neighborhood sorting and predicts fold-wise global means directly in $O(K n q)$ time.
	
	\item \textbf{Granular-Ball Dependency Evaluation ($T_\Gamma(s)$):}
	Evaluating $\Gamma(P)$ for a candidate $P$ begins with aggregating the symmetric directed dissimilarity matrix $d_P$, taking $O(n^2 s)$ operations. At a granular node containing $b$ instances, medoid discovery and farthest-seed selection require $O(b^2)$ distance lookups, while dispersion updates require $O(b q)$. While the minimum support constraint $s_{\min} \ge 2$ bounds the total number of binary splits by $\lfloor n/s_{\min} \rfloor - 1$, in the worst-case scenario of a completely degenerate, unbalanced tree, summing over all hierarchical levels yields:
	\begin{equation}
		T_\Gamma(s) = O\left(n^2 s + n^3 + n^2 q\right).
		\label{eq:eval-gamma-complexity}
	\end{equation}
	Under balanced partitioning, this bound naturally tightens toward $O(n^2 s + n^2 \log(n/s_{\min}) + n q \log(n/s_{\min}))$.
\end{enumerate}

\subsubsection{Combinatorial Search Procedures}
Equipped with the single-query oracles $T_R = T_R(m)$ and $T_\Gamma = T_\Gamma(m)$, we characterize the overall time and memory requirements of the reduction algorithms.

\begin{proposition}[Algorithmic Complexity Bounds]
	\label{prop:complexity-bounds}
	The asymptotic time complexity and auxiliary memory requirements of the proposed reduction algorithms are characterized as follows:
\begin{enumerate}[label=(\roman*)]
		\item \textbf{Dependency-Guided Heuristic:} 
		The forward expansion evaluates at most $\frac{1}{2}m(m+1)$ candidate additions. The subsequent restarting backward pruning commits at most $m$ successful deletions alongside a final scan of at most $m$ checks, resulting in $O(m^2)$ total subset queries. Evaluating dependency $\Gamma$ strictly on loss-improving candidates ($J(P)$) yields an overall worst-case execution time of:
		\begin{equation}
			T_{\mathrm{guided}} = T_{\mathrm{prep}} + O\left(m^2 (T_R + T_\Gamma)\right).
			\label{eq:guided-complexity}
		\end{equation}
		Maintaining subset caches requires $O(m^3)$ scalar indexing storage, ensuring a total memory bound of $O(m n^2 + n q + n m + m^3)$.
		
		\item \textbf{Exact Branch-and-Bound:} 
		The binary inclusion/exclusion tree comprises at most $2^{m+1}-1$ states. Because feasibility depends strictly on $R_{\mathrm{inner}}$, the exact search entirely bypasses the calculation of $\Gamma(P)$. Pruning nodes exceeding the incumbent cost $C^*$ substantially reduces the practical search frontier, while maintaining the theoretical worst-case bound:
		\begin{equation}
			T_{\mathrm{exact}} = T_{\mathrm{prep}} + O\left(2^m (T_R + m)\right),
			\label{eq:exact-complexity}
		\end{equation}
		with depth-first execution requiring $O(m^2)$ recursion stack memory when executed without a global exponential subset cache.
		
		\item \textbf{Budgeted Exchange Refinement:} 
		For a subset of size $s$, the 1-exchange and deletion neighborhood contains $|\mathcal{M}| \le s + s(m - s) = O(m^2)$ candidate transformations. Gated by an explicit maximum evaluation budget $B_{\mathrm{eval}}$, the refinement procedure strictly executes within:
		\begin{equation}
			T_{\mathrm{refine}} = O\left(B_{\mathrm{eval}} T_R + B_{\mathrm{eval}} m^2 \log m\right),
			\label{eq:refine-complexity}
		\end{equation}
		with auxiliary candidate list storage bounded by $O(m^2)$.
	\end{enumerate}
\end{proposition}

\subsubsection{Theoretical Boundaries and Practical Considerations}
The complexity analysis highlights a clear decoupling between operational acquisition costs $C(P)$ and computational running times:
\begin{itemize}
	\item \textbf{Scalability with Data Volume ($n$):} The offline pairwise dissimilarity caching strategy effectively trades memory for time, amortizing the $O(m n^2)$ distance construction over all subsequent subset traversals. This ensures that candidate evaluation reduces to lightweight index manipulation and matrix lookups.
	\item \textbf{Scalability with Dimensionality ($m$):} For high-dimensional datasets ($m \gg 1$), the quadratic query bound $O(m^2)$ of Algorithm~\ref{alg2:guided} provides a computationally tractability guarantee compared to exponential complete solvers, while the dependency oracle $\Gamma$ focuses search trajectories on distributionally discriminative sub-manifolds without introducing extra tuning hyperparameters.
\end{itemize}
